\documentclass[11pt,compress,t,notes=noshow, xcolor=table]{beamer}

\input{../../style/preamble}
\input{../../latex-math/basic-math}
\input{../../latex-math/optim-basics}

\title{Optimization in Machine Learning}

\begin{document}

\titlemeta{
Foundations
}{
Quadratic functions
}{
figure/quadratic_functions_2D_example_1_1.png
}{
\item Definition of quadratic functions
\item Gradient, Hessian
\item Optima
\item Geometry: spectrum and curvature
\item Condition number
}

% ==========================================================================
% Part I: definition, basic properties, optima
% ==========================================================================

\begin{framei}{Univariate quadratic}
\item Quadratic function $q: \R \to \R$
$$ q(x) = a x^2 + b x + c,\quad a \ne 0 $$
\imageC[0.35]{figure/quadratic_functions_1D.png}
\item Left: $q_1(x) = x^2$. Right: $q_2(x) = - x^2$
\end{framei}


\begin{framei}{Univariate: basic properties}
\item Slope at $(x, q(x))$: $$ q'(x) = 2 a x + b $$
\imageC[0.25]{figure/quadratic_functions_1D_derivative.png}
\item Curvature: $$ q''(x) = 2 a $$
\imageC[0.25]{figure/quadratic_functions_1D_curvature.png}
\item $a > 0$: $q$ convex, bounded from below, unique global minimum
\item $a < 0$: $q$ concave, bounded from above, unique global maximum
\item Optimum $\xs$
$$ q'(x) = 0 \;\Leftrightarrow\; 2 a x + b = 0 \;\Rightarrow\; \xs = \frac{-b}{2 a} $$\\
as 2nd derivative: $q''(\xs) = 2 a \ne 0$
\end{framei}


\begin{framei}{Multivariate quadratic}
\item $q: \R^d \to \R$
$$ q(\xv) = \xv^T \A \xv + \bv^T \xv + c $$
\item with $\A \in \R^{d \times d}$ full rank, $\bv \in \R^d$, $c \in \R$
\vfill
\splitV{
\imageC[0.9]{figure/quadratic_functions_2D_example_1_1.png}
}{
\imageC[0.9]{figure/quadratic_functions_2D_example_1_2.png}
}
\end{framei}


\begin{framei}{Symmetrization}
\item W.l.o.g. assume $\A$ symmetric, i.e., $\A^T = \A$
\item If $\A$ not symmetric, there exists symmetric $\tilde \A$ with $$ q(\xv) = \xv^T \A \xv = \xv^T \tilde \A \xv =: \tilde q(\xv) $$
\item Justification
$$ q(\xv) = \xv^T \A \xv = \tfrac12 \xv^T \underbrace{(\A + \A^T)}_{\tilde \A_1} \xv + \tfrac12 \xv^T \underbrace{(\A - \A^T)}_{\tilde \A_2} \xv $$
\item $\tilde \A_1$ symmetric, $\tilde \A_2$ anti-symmetric ($\tilde \A_2^T = - \tilde \A_2$)
\item Since $\xv^T \A^T \xv$ is a scalar, equal to its transpose
\begin{align*}
\xv^T (\A - \A^T) \xv &= \xv^T \A \xv - \xv^T \A^T \xv = \xv^T \A \xv - (\xv^T \A^T \xv)^T \\
                      &= \xv^T \A \xv - \xv^T \A \xv = 0
\end{align*}
\item Therefore $\tilde q(\xv) = \xv^T \tilde \A \xv$ with $\tilde \A = \tilde \A_1/2$
\end{framei}


\begin{framei}{Gradient and Hessian}
\item $q: \R^d \to \R$
$$ q(\xv) = \xv^T \A \xv + \bv^T \xv + c $$
\item Gradient
$$ \nabla q(\xv) = ((\A^T + \A) \xv + \bv)^T $$
\item Under assumed symmetry: $\nabla q(\xv) = (2 \A \xv + \bv)^T$
\item Directional derivative: $\nabla q(\xv)\, \vv $
\item Hessian
$$ \nabla^2 q(\xv) = (\A^T + \A) = 2 \A =: \H \in \R^{d \times d} $$
\item Under assumed symmetry: $\H = 2 \A$
\item Directional curvature: $\vv^T \H \vv $
\end{framei}


\begin{framei}{Optimum}
\item $q: \R^d \to \R$
$$ q(\xv) = \xv^T \A \xv + \bv^T \xv + c $$
\item Since $\A$ full rank, unique stationary point $\xvs$ (min, max, or saddle)
\begin{align*}
\nabla q(\xvs)      &= \boldsymbol{0}^T \\
(2 \A \xvs + \bv)^T &= \boldsymbol{0}^T \\
\xvs                &= -\tfrac{1}{2} \A^{-1} \bv
\end{align*}
\item $q(\xvs) = c - \frac12 \bv^T \A^{-1} \bv$
\vfill
\imageC[1]{figure_man/minmaxsaddle.png}
\item Left: $\A$ pos. def. \quad Middle: $\A$ neg. def. \quad Right: $\A$ indefinite
\end{framei}


\begin{framei}[fs=footnotesize]{Optima: rank-deficient case}
\item $q: \R^d \to \R$
$$ q(\xv) = \xv^T \A \xv + \bv^T \xv + c $$
\item Assume $\A$ symmetric now
\item For stationary points to exist, we need : $\nabla q(\xv) = 2 \A \xv + \bv = 0$
\item This implies we need $\bv \in range(\A)$, let's assume this is the case
\item Let $\xv_p$ be stationary, so $2 \A \xv_p = -\bv$
\item Then any point in affine space $\xv_p + ker(\A)$ is also stationary, with same function value and same Hessian (as it is constant)
\imageC[0.35]{figure_man/convex-example.png}
\end{framei}


\begin{framei}[fs=footnotesize]{Optima: rank-deficient case}
\item All affine spaces of form $\xv_p + ker(\A)$ for diff. valid $\xv_p$ are the same
\item Any stationary point must be in $\xv_p + ker(\A)$
\item So $\xv_p + ker(\A)$ are all the stationary points, with same curvature
\item If $\A \succeq 0$, they are all minima
\item If $\A \preceq 0$, they are all maxima
\item If $\A$ is indefinite, they are all saddle points
\imageC[0.35]{figure_man/convex-example.png}
\end{framei}


% ==========================================================================
% Part II: geometry, spectrum, condition
% ==========================================================================

\begin{framei}{Geometry of quadratic functions}
\item Example: $\Amat = \begin{pmatrix} 2 & -1 \\ -1 & 2\end{pmatrix}$ $\Rightarrow$ $\H = 2\Amat = \begin{pmatrix} 4 & -2 \\ -2 & 4\end{pmatrix}$
\item Since $\H$ symmetric: eigendecomposition $\H = \V\Lambda\V^T$
\begin{equation*}
\V = \begin{pmatrix}
| & | \\
\textcolor{magenta}{v_\text{max}} & \textcolor{orange}{v_\text{min}} \\
| & |
\end{pmatrix}
= \frac{1}{\sqrt{2}} \begin{pmatrix}
1 & 1 \\
-1 & 1
\end{pmatrix}\;\text{orthogonal}
\end{equation*}
\begin{equation*}
\Lambda = \begin{pmatrix}
\textcolor{magenta}{\lambda_\text{max}} & 0 \\
0 & \textcolor{orange}{\lambda_\text{min}}
\end{pmatrix}
= \begin{pmatrix}6 & 0 \\ 0 & 2\end{pmatrix}
\end{equation*}
\imageC[0.35]{figure/quadr-eigenv.png}
\end{framei}


\begin{framei}[fs=footnotesize]{Spectrum and curvature}
\item $\vvmax$ direction of highest curvature, with curvature value $\lammax$
$$
\vv^T \H \vv = \vv^T \V \Lambda \V^T \vv = \wv^T \Lambda \wv =
\sumid \lambda_i w_i^2 \le \lammax \sumid w_i^2 = \lammax || \wv ||^2
$$
\item Since $|| \vv || = || \xv ||$ ($\V$ orthogonal):
$\max_{|| \vv || = 1} \vv^T \H \vv \le \lammax$
\item For $\vvmax$ we obtain this upper bound:
$\vvmax^T \H \vvmax = \mathbf{e}_1^T \Lambda \mathbf{e}_1 = \lammax$
\item Analogously, $\vvmin$ direction of lowest curvature, with curvature value $\lammin$
\item Contour lines of any quadratic function are ellipses
\end{framei}


\begin{framei}[fs=footnotesize]{Second-order sufficient condition}
\item Recall: Second order condition for optimality is sufficient
\item If $H(\xv^\ast) \succ 0$ at stationary point $\xv^\ast$, then $\xv^\ast$ local minimum ($\prec$ for maximum)
\item[] \begin{equation*}
f(\xv) = f(\xv^\ast) + \underbrace{\nabla f(\xv^\ast)}_{=0}(\xv-\xv^\ast) + \tfrac{1}{2}\underbrace{(\xv-\xv^\ast)^T H(\xv^\ast)(\xv-\xv^\ast)}_{\ge \textcolor{orange}{\lambda_\text{min}} \|\xv-\xv^\ast\|^2} + \underbrace{R_2(\xv,\xv^\ast)}_{=o(\|\xv-\xv^\ast\|^2)}
\end{equation*}
\item Choose $\eps>0$ s.t. $|R_2(\xv,\xv^\ast)| < \tfrac{1}{2} \textcolor{orange}{\lambda_\text{min}} \|\xv-\xv^\ast\|^2$ for $\xv \neq \xv^\ast$, $\|\xv-\xv^\ast\|<\eps$
\item[] \begin{equation*}
f(\xv) \ge f(\xv^\ast) + \underbrace{\tfrac{1}{2} \textcolor{orange}{\lambda_\text{min}} \|\xv-\xv^\ast\|^2 + R_2(\xv,\xv^\ast)}_{>0} > f(\xv^\ast)
\end{equation*}
\end{framei}


\begin{framei}{Eigenvalues and shape}
\item If spectrum of $\Amat$ is known, also that of $\H = 2\Amat$ is known
\item If all eigenvalues of $\H$ $\overset{(>)}{\ge 0}$ ($\Leftrightarrow$ $\H \overset{(\succ)}{\succcurlyeq} 0$):
\begin{itemize}
\item $q$ (strictly) convex
\item (Unique) global minimum
\end{itemize}
\item If all eigenvalues of $\H$ $\overset{(<)}{\le 0}$ ($\Leftrightarrow$ $\H \overset{(\prec)}{\preccurlyeq} 0$):
\begin{itemize}
\item $q$ (strictly) concave
\item (Unique) global maximum
\end{itemize}
\item If $\H$ has both positive and negative eigenvalues ($\Leftrightarrow$ $\H$ indefinite):
\begin{itemize}
\item $q$ neither convex nor concave
\item there is a saddle point
\end{itemize}
\imageC[0.67]{figure_man/minmaxsaddle.png}
\end{framei}


\begin{framei}{Condition and curvature}
\item $\kappa(\H) = \kappa(\Amat) = |\textcolor{magenta}{\lambda_\text{max}}| / |\textcolor{orange}{\lambda_\text{min}}|$
\item High condition means
\begin{itemize}
\item $|\textcolor{magenta}{\lambda_\text{max}}| \gg |\textcolor{orange}{\lambda_\text{min}}|$
\item Curvature along $\textcolor{magenta}{\vv_\text{max}} \gg$ along $\textcolor{orange}{\vv_\text{min}}$
\item Problem for algorithms like gradient descent
\end{itemize}
\imageC[1]{figure_man/quadr-conds.png}
\begin{center}\begin{footnotesize} Left: Excellent condition. Middle: Good condition. Right: Bad condition.\end{footnotesize}\end{center}
\end{framei}


\begin{framei}{Approximation of smooth functions}
\item Any $f \in \mathcal{C}^2$ can be locally approximated by quadratic function (second order Taylor)
\item[] \begin{equation*}
f(\xv) \approx f(\bm{\tilde{x}}) + \nabla f(\bm{\tilde{x}})(\xv-\bm{\tilde{x}}) + \tfrac 12(\xv-\bm{\tilde{x}})^T\nabla^2 f(\bm{\tilde{x}})(\xv-\bm{\tilde{x}})
\end{equation*}
\imageC[0.3][https://daniloroccatano.blog]{figure_man/taylor_2D_quadratic.png}
\begin{center}\begin{footnotesize} $f$ and second order approximation: dark vs bright grid.\end{footnotesize}\end{center}
\item $\implies$ Hessians provide information about \textbf{local} geometry of a function
\end{framei}

\endlecture
\end{document}
